\documentclass[11pt]{article} % Documento de tipo "artículo", con fuente por defecto de 11pt

\usepackage[utf8]{inputenc}   % Permite el uso de caracteres UTF8
\usepackage[T1]{fontenc}      % Usa fuentes "ricas", que facilitan el copy-paste en la salida
\usepackage[spanish]{babel}   % Traducción al español de meses, funciones matemáticas, etc.
\usepackage{amsmath}          % Herramientas matemáticas extra, como los entornos "align" y "gather"
\usepackage{graphicx}         % Permite la inserción de imágenes
\usepackage{hyperref}         % Convierte las referencias en hiperenlaces y añade el comando \url
\usepackage{xcolor}           % Permite cambiar el color del texto
\usepackage{tgpagella}        % Fuente TeX Gyre Pagella, usada en todo el documento
\usepackage{fancyhdr}         % Para personalizar las cabeceras y el pie de página
\usepackage[
  headheight=3cm,
  top=4.5cm,
  headsep=1.5cm,
  bottom=3cm
]{geometry}                   % Establece márgenes personalizados

% Colores definidos en la guía de estilo del IEEE
% (https://brand-experience.ieee.org/guidelines/digital/style-guide/branding-visual-elements/)
\definecolor{ieee-blue}{RGB}{0, 98, 155}
\definecolor{ieee-gray}{RGB}{117, 120, 123}

% Título y fecha del documento
\newcommand{\theplace}{Madrid}
\newcommand{\thedate}{19 de Septiembre de 2018}
\newcommand{\thetitle}{Recuperación Cuenta de Twitter}

% Ajustes de la cabecera y el pie de página
\pagestyle{fancy}
\fancyhf{} % Elimina estilos de pie y cabecera por defecto
\lhead{\includegraphics[height=40pt]{./LogoIEEEsb.eps}}
\rhead{
  \color{ieee-blue}
  Rama Universitaria del IEEE de Madrid
  \\[4pt]
  \color{ieee-gray}
  \textit{\thetitle} \\
  \theplace, \thedate
}
\cfoot{
  \scriptsize
  \textsc{Rama Universitaria del IEEE de Madrid} \\
  Escuela Técnica Superior de Ingenieros de Telecomunicación - Universidad Politécnica de Madrid \\
  Av. Complutense, 30 (A101.4-L), 28040 Madrid / 91 067 2681 / ieeesb.es}
\renewcommand{\headrulewidth}{0pt} % Quita la línea por defecto debajo de la cabecera


\begin{document}
\begin{center}
  \bf \LARGE
  \MakeUppercase{\thetitle}
\end{center}
    Carlos Cuenca Enrique, Presidente de la Rama Universitaria del IEEE de Madrid,es autorizado por la organización para el uso y administración de la cuenta \textbf{@ieeesb}.\\ \\
    En el momento de la creación de la cuenta, la organización Rama Universitaria IEEE de Madrid autorizó la creación de la misma y su posterior uso.
    La organización ha intentado acceder de diversos medios a la cuenta y ha abierto un ticket de soporte desde la cuenta oficial de la organización para estos asuntos. Este documento es el resultado del proceso de recuperación de la cuenta iniciado por la organización. \\ \\
    Solicitamos que la dirección de correo electrónico asociada a la cuenta \textbf{@ieeesb}  sea sustituida por la que nos hemos puesto en contacto con ustedes para este proceso: \\ \textbf{cuentas@ieeesb.es}. \\ \\
    
    La organización: \hspace{200pt} Firmado: \\\\\\\\\\\\\\\ Rama Universitaria del IEEE \hspace{150pt}Carlos Cuenca Enrique
\end{document}